% sql4-vs-sqlite.tex
%
% SQL4 vs. SQLite: a scaling comparison
% Fritz Henglein, Department of Computer Science, University of Copenhagen (DIKU)
% Advanced Programming, DIKU, 2026-10-20
%
% Standalone slide deck, separate from generic-discrimination.tex, reusing the
% same KU look and feel. Shows the correspondence between SQL queries and
% SQL4 (src/SQL4.hs) functions, and the measured run times of SQL4 (compiled
% with -O2) against SQLite (data bulk-loaded via its fastest CSV .import path
% under synchronous=OFF/journal_mode=OFF/temp_store=MEMORY) on identical
% generated inputs.
% Compile with:   pdflatex sql4-vs-sqlite.tex

\documentclass{beamer}

% ---------------------------------------------------------------------------
% KU look and feel (same as generic-discrimination.tex)
% ---------------------------------------------------------------------------
\useinnertheme[shadow=true]{rounded}

\definecolor{kugreen}{RGB}{50,93,61}
\usecolortheme[named=kugreen]{structure}
\setbeamercolor{frametitle}{fg=kugreen}
\setbeamercolor{title}{fg=white,bg=kugreen}
\setbeamercolor{block title}{fg=white,bg=kugreen}
\setbeamercolor{block body}{fg=black,bg=kugreen!25}

\setbeamerfont{frametitle}{size=\Large,series=\bfseries}
\setbeamerfont{block title}{size=\normalsize}
\setbeamertemplate{navigation symbols}{}

\usepackage{listings}
\lstset{%
  language=Haskell,
  keywords={case,class,data,deriving,do,else,if,import,in,infix,infixl,
            infixr,instance,let,module,newtype,of,then,type,where},
  basicstyle=\upshape\ttfamily,
  keywordstyle=\color{kugreen},
  columns=fixed,
  basewidth=0.5em,
  keepspaces=true,
  showstringspaces=false}

\usepackage{booktabs}
\usepackage{pgfplots}
\pgfplotsset{compat=1.18}

\newcommand{\kubackground}[1]{%
  \setbeamertemplate{background}{%
    \includegraphics[width=\paperwidth,height=\paperheight]{#1}}}
\kubackground{figures/ku-background}

\newlength{\kutitletop}     \setlength{\kutitletop}{16.54pt}
\newlength{\kutitleindent}  \setlength{\kutitleindent}{14.54pt}
\newlength{\kutitlebottom}  \setlength{\kutitlebottom}{-6.61pt}
\setbeamertemplate{frametitle}{%
  \vskip\kutitletop
  \begin{beamercolorbox}[wd=\paperwidth,leftskip=\kutitleindent,
                         rightskip=0pt plus 1fil]{frametitle}%
    \usebeamerfont{frametitle}\strut\insertframetitle\strut\par
  \end{beamercolorbox}%
  \vskip\kutitlebottom}

\newlength{\kufootheight}   \setlength{\kufootheight}{13.92pt}
\setbeamertemplate{footline}{%
  \vbox to\kufootheight{%
    \vfil
    \hbox to\paperwidth{\hskip5pt\usebeamerfont{footline}\tiny
      \insertframenumber\hfil}%
    \vskip10pt}}

\newlength{\kutitlepagetop}  \setlength{\kutitlepagetop}{39.45pt}
\newlength{\kudateskip}      \setlength{\kudateskip}{2pt}
\setbeamertemplate{title page}{%
  \vskip\kutitlepagetop
  \begingroup
    \centering
    \begin{beamercolorbox}[sep=4pt,center,colsep=-4bp,rounded=true,
                           shadow=true]{title}
      \usebeamerfont{title}\inserttitle\par
    \end{beamercolorbox}%
    \vskip1em\par
    \begin{beamercolorbox}[sep=8pt,center]{author}
      \usebeamerfont{author}\insertauthor
    \end{beamercolorbox}
    \begin{beamercolorbox}[sep=8pt,center]{institute}
      \usebeamerfont{institute}\insertinstitute
    \end{beamercolorbox}
    \vskip\kudateskip
    \begin{beamercolorbox}[sep=8pt,center]{date}
      \usebeamerfont{date}\insertdate
    \end{beamercolorbox}
  \endgroup
  \vskip0pt plus 1filll}

% ---------------------------------------------------------------------------
% Shared style for the three scaling plots: linear-linear axes, grouped bars
% (SQLite query, SQLite load, SQL4), one bar triple per measured input size n.
% ---------------------------------------------------------------------------
\pgfplotsset{
  benchplot/.style={
    width=\textwidth, height=0.72\textheight,
    xlabel={$n$ (input size, millions)},
    ylabel={time (s)},
    ybar,
    bar width=4pt,
    enlarge x limits={abs=10pt},
    ymajorgrids,
    axis x line*=bottom,
    axis y line*=left,
    legend style={at={(0.02,0.98)},anchor=north west,font=\scriptsize,
                  draw=kugreen,fill=white},
    tick label style={font=\scriptsize},
    label style={font=\scriptsize},
  }
}

\title{SQL4 vs.\ SQLite}
\author{Fritz Henglein\\ \texttt{henglein@diku.dk}}
\institute{Department of Computer Science\\ University of Copenhagen\\ (DIKU)}
\date{Advanced Programming, DIKU, 2026-10-20}

\begin{document}

% ===========================================================================
{\kubackground{figures/ku-background-title}
\begin{frame}[plain]
  \titlepage
\end{frame}
}

% ===========================================================================
\begin{frame}[fragile]{SQL queries and SQL4 functions}
  Rows: $(\mathit{key}, \mathit{payload})$ pairs with $\mathit{key}$ uniform
  on $[0, 1000)$;\\
  \lstinline|eKey = NatE 999|\quad (equivalence on keys)\\
  \lstinline|eRow = MapE fst (NatE 999)|\quad (equivalence on whole rows)
  \vspace{0.3cm}

  \renewcommand{\arraystretch}{1.6}
  \begin{center}
  \scriptsize
  \begin{tabular}{@{}p{0.37\textwidth}@{\hspace{0.8em}}p{0.57\textwidth}@{}}
    \toprule
    \textbf{SQL query} & \textbf{SQL4 (\texttt{src/SQL4.hs}) function} \\
    \midrule
    \texttt{SELECT * FROM r, s}\newline
    \texttt{WHERE r.key = s.key} &
    \lstinline|select ((fst,fst) `Is` eKey) (r `X` s)|\newline
    \emph{(dispatches to \lstinline|djoin|)} \\
    \texttt{SELECT key, COUNT(*)}\newline
    \texttt{FROM r GROUP BY key} &
    \lstinline|perform (Func count) (groupBy eRow r)| \\
    \texttt{SELECT DISTINCT key FROM r} &
    \lstinline|coalesceBy eKey (perform (Func fst) r)| \\
    \bottomrule
  \end{tabular}
  \end{center}
\end{frame}

% ===========================================================================
\begin{frame}{Performance: equi-join}
  \begin{center}
  \begin{tikzpicture}
    \begin{axis}[benchplot,
        xtick={0,4,8,12,16,20,24},
      ]
      \addplot+[bar shift=-5pt, fill=kugreen!70, draw=kugreen] coordinates {
        (0.4,0.064) (0.8,0.128) (1.6,0.284) (3.2,0.537)
        (6.4,1.227) (12.8,2.885) (23,5.978)
      };
      \addplot+[bar shift=0pt, fill=kugreen!25, draw=kugreen] coordinates {
        (0.4,0.340) (0.8,0.666) (1.6,1.348) (3.2,2.703)
        (6.4,5.489) (12.8,10.852) (23,19.893)
      };
      \addplot+[bar shift=5pt, fill=gray!60, draw=black] coordinates {
        (0.4,0.170) (0.8,0.348) (1.6,0.734) (3.2,1.517)
        (6.4,3.278) (12.8,7.082) (23,14.475)
      };
      \legend{SQLite query, SQLite load, SQL4}
    \end{axis}
  \end{tikzpicture}
  \end{center}
  \footnotesize \texttt{SELECT * FROM r, s WHERE r.key = s.key}
  (cardinality only: enumerating the result is intractable, e.g.\ $5.29
  \times 10^{11}$ matches at $n = 23{,}000{,}000$)
\end{frame}

% ===========================================================================
\begin{frame}{Performance: group by + count}
  \begin{center}
  \begin{tikzpicture}
    \begin{axis}[benchplot,
        xtick={0,10,20,30,40,50,61},
      ]
      \addplot+[bar shift=-5pt, fill=kugreen!70, draw=kugreen] coordinates {
        (1,0.148) (4,0.691) (16,3.629) (32,8.671)
        (48,13.489) (61,17.528)
      };
      \addplot+[bar shift=0pt, fill=kugreen!25, draw=kugreen] coordinates {
        (1,0.840) (4,3.370) (16,13.800) (32,28.079)
        (48,42.812) (61,55.322)
      };
      \addplot+[bar shift=5pt, fill=gray!60, draw=black] coordinates {
        (1,0.149) (4,0.673) (16,3.279) (32,7.050)
        (48,11.001) (61,18.233)
      };
      \legend{SQLite query, SQLite load, SQL4}
    \end{axis}
  \end{tikzpicture}
  \end{center}
  \footnotesize \texttt{SELECT key, COUNT(*) FROM r GROUP BY key}
\end{frame}

% ===========================================================================
\begin{frame}{Performance: distinct}
  \begin{center}
  \begin{tikzpicture}
    \begin{axis}[benchplot,
        xtick={0,50,100,150,195},
      ]
      \addplot+[bar shift=-5pt, fill=kugreen!70, draw=kugreen] coordinates {
        (1,0.118) (4,0.437) (16,1.754) (64,6.660)
        (128,13.976) (195,23.704)
      };
      \addplot+[bar shift=0pt, fill=kugreen!25, draw=kugreen] coordinates {
        (1,0.840) (4,3.370) (16,13.800) (64,57.918)
        (128,114.855) (195,178.170)
      };
      \addplot+[bar shift=5pt, fill=gray!60, draw=black] coordinates {
        (1,0.087) (4,0.347) (16,1.387) (64,5.486)
        (128,11.549) (195,21.116)
      };
      \legend{SQLite query, SQLite load, SQL4}
    \end{axis}
  \end{tikzpicture}
  \end{center}
  \footnotesize \texttt{SELECT DISTINCT key FROM r}
\end{frame}

\end{document}
